% !TEX program = xelatex
%==============================================================================
%  Phân công code vấn đáp đồ án X509Certificate
%  Phần việc chi tiết: NGUYỄN LÊ ANH KIỆT
%  — Vòng đời chứng chỉ, Thu hồi & Hạ tầng CRL/OCSP —
%
%  Biên dịch:  xelatex PhanCong_NguyenLeAnhKiet.tex  (chạy 2 lần để có mục lục)
%==============================================================================
\documentclass[12pt,a4paper]{article}

% ---- Font (XeLaTeX, hỗ trợ Unicode tiếng Việt đầy đủ) -----------------------
\usepackage{fontspec}
\setmainfont{Latin Modern Roman}
\setsansfont{Latin Modern Sans}
\setmonofont{Latin Modern Mono}

% ---- Bố cục trang -----------------------------------------------------------
\usepackage[a4paper,margin=2.4cm]{geometry}
\usepackage{setspace}\setstretch{1.18}
\usepackage{parskip}
\usepackage{amssymb} % \square cho checklist

% ---- Màu sắc ----------------------------------------------------------------
\usepackage{xcolor}
\definecolor{accent}{HTML}{1F4E79}
\definecolor{accentlt}{HTML}{2E75B6}
\definecolor{rulegray}{HTML}{C3CEDA}
\definecolor{codebg}{HTML}{F6F8FA}
\definecolor{codeframe}{HTML}{D9E1EA}
\definecolor{cmt}{HTML}{6A737D}
\definecolor{kw}{HTML}{0B5394}
\definecolor{str}{HTML}{A31515}
\definecolor{num}{HTML}{098658}
\definecolor{noteblue}{HTML}{EAF2FB}
\definecolor{qabg}{HTML}{F3EEFA}
\definecolor{qarule}{HTML}{7E57C2}
\definecolor{warnbg}{HTML}{FBF1E6}
\definecolor{warnrule}{HTML}{D9822B}
\definecolor{okbg}{HTML}{E8F5E9}
\definecolor{okrule}{HTML}{2E7D32}

% ---- Tiêu đề mục ------------------------------------------------------------
\usepackage{titlesec}
\titleformat{\section}{\Large\bfseries\color{accent}}{\thesection}{0.6em}{}
\titleformat{\subsection}{\large\bfseries\color{accentlt}}{\thesubsection}{0.6em}{}
\titleformat{\subsubsection}{\normalsize\bfseries\color{accentlt}}{\thesubsubsection}{0.5em}{}
\titlespacing*{\section}{0pt}{1.4em}{0.6em}

% ---- Danh sách --------------------------------------------------------------
\usepackage{enumitem}
\setlist[itemize]{leftmargin=1.4em,topsep=2pt,itemsep=2pt}
\setlist[enumerate]{leftmargin=1.6em,topsep=2pt,itemsep=2pt}

% ---- Bảng -------------------------------------------------------------------
\usepackage{booktabs}
\usepackage{tabularx}
\usepackage{array}
\usepackage{colortbl}
\newcolumntype{L}[1]{>{\raggedright\arraybackslash}p{#1}}

% ---- Hộp callout (tcolorbox) ------------------------------------------------
\usepackage[most]{tcolorbox}
\newtcolorbox{notebox}{colback=noteblue,colframe=accentlt,boxrule=0pt,
  leftrule=3pt,arc=2pt,left=8pt,right=8pt,top=5pt,bottom=5pt}
\newtcolorbox{qabox}{colback=qabg,colframe=qarule,boxrule=0pt,
  leftrule=3pt,arc=2pt,left=8pt,right=8pt,top=5pt,bottom=5pt}
\newtcolorbox{warnbox}{colback=warnbg,colframe=warnrule,boxrule=0pt,
  leftrule=3pt,arc=2pt,left=8pt,right=8pt,top=5pt,bottom=5pt}

% ---- Mã nguồn (fvextra Verbatim — render đúng tiếng Việt dưới XeLaTeX) -------
% Lưu ý: listings reorder/garble dấu tiếng Việt trong code, nên dùng fvextra.
\usepackage{fvextra}
\fvset{
  numbers=left,
  numbersep=8pt,
  frame=single,
  framesep=6pt,
  rulecolor=\color{codeframe},
  fontsize=\footnotesize,
  breaklines=true,
  breakanywhere=true,
  breaksymbolleft={\footnotesize\color{cmt}\ensuremath{\hookrightarrow}},
}
% \code: định danh inline; underscore tự cho phép ngắt dòng (tránh tràn bảng).
\newcommand{\code}[1]{{\ttfamily\small\def\_{\textunderscore\discretionary{}{}{}}#1}}

% ---- Liên kết ---------------------------------------------------------------
\usepackage{hyperref}
\hypersetup{
  colorlinks=true, linkcolor=accent, urlcolor=accentlt, citecolor=accent,
  pdftitle={Phan cong code van dap - Nguyen Le Anh Kiet - X509Certificate},
  pdfauthor={Nguyen Le Anh Kiet},
}

% ---- Header / footer --------------------------------------------------------
\usepackage{fancyhdr}
\pagestyle{fancy}\fancyhf{}
\renewcommand{\headrulewidth}{0.4pt}
\renewcommand{\footrulewidth}{0pt}
\fancyhead[L]{\footnotesize\color{accent}X509Certificate — Phân công vấn đáp}
\fancyhead[R]{\footnotesize\color{accent}Nguyễn Lê Anh Kiệt}
\fancyfoot[C]{\thepage}

% ---- Hộp nhỏ dùng cho sơ đồ (thay cho TikZ, biên dịch ổn định) --------------
\newcommand{\flowbox}[2]{\colorbox{#1}{\parbox{8.4cm}{\centering #2}}}
\newcommand{\stepbox}[1]{\fcolorbox{accent}{codebg}{\footnotesize\strut #1}}

\begin{document}

%==============================================================================
% TRANG BÌA
%==============================================================================
\begin{titlepage}
\centering
\vspace*{2.2cm}
{\color{accent}\rule{\linewidth}{1.2pt}}\\[0.6cm]
{\Huge\bfseries Phân công code vấn đáp đồ án}\\[0.35cm]
{\Huge\bfseries\color{accent} X509Certificate}\\[0.5cm]
{\Large Phần việc chi tiết của thành viên}\\[0.3cm]
{\LARGE\bfseries Nguyễn Lê Anh Kiệt}\\[0.45cm]
{\large\itshape Vòng đời chứng chỉ, Thu hồi\\ và Hạ tầng CRL \& OCSP}\\[0.5cm]
{\color{accent}\rule{\linewidth}{1.2pt}}\\[1.4cm]

\begin{minipage}{0.85\linewidth}
\centering\small
Tài liệu này mô tả chi tiết phần code, lý thuyết X.509 cần nắm, các test case
và cách tích hợp cho phần việc \textbf{vòng đời chứng chỉ (lifecycle), thu hồi
(revocation) và hạ tầng phân phối trạng thái CRL \& OCSP}, gắn trực tiếp với mã
nguồn thực tế của đồ án trong các thư mục \code{src/core/} (\code{crl.py}),
\code{src/services/} (\code{cert\_lifecycle.py}, \code{revocation\_workflow.py},
\code{crl\_publish.py}, \code{infra\_manager.py}), \code{src/infra/}
(\code{crl\_server.py}, \code{ocsp\_server.py}) và các màn hình UI tương ứng.
\end{minipage}

\vfill
{\large Môn: Mã hoá và Ứng dụng (MHUD)}\\[0.2cm]
{\large Học phần đồ án X.509 Certificate}\\[0.6cm]
{\small Ngày biên soạn: 14/06/2026}
\end{titlepage}

\tableofcontents
\thispagestyle{fancy}
\clearpage

%==============================================================================
\section{Mục tiêu và phạm vi phần việc}
%==============================================================================

Theo bảng phân công chung, thành viên \textbf{Nguyễn Lê Anh Kiệt} phụ trách
cụm chức năng \textbf{vòng đời chứng chỉ, thu hồi và hạ tầng CRL \& OCSP}. Đây
là phần \emph{nằm ở khâu hậu cấp phát}: sau khi một chứng chỉ đã được Root CA
ký và phát hành cho khách hàng, nó vẫn cần được quản lý suốt vòng đời còn lại
--- gia hạn (renew), thu hồi (revoke) khi khoá bị lộ hay không còn dùng nữa, và
quan trọng nhất là \emph{công bố trạng thái thu hồi ra ngoài} để bất kỳ bên nào
cũng kiểm tra được một cert còn hợp lệ hay không.

Cụm này trả lời câu hỏi cốt lõi của PKI: \emph{``Làm sao biết một chứng chỉ tuy
chưa hết hạn nhưng đã bị thu hồi?''}. Câu trả lời gồm hai cơ chế bù nhau:
\textbf{CRL} (Certificate Revocation List --- danh sách thu hồi do CA ký, phát
hành định kỳ) và \textbf{OCSP} (Online Certificate Status Protocol --- tra cứu
trạng thái từng serial theo thời gian thực). Phần việc của Kiệt hiện thực trọn
vẹn cả hai từ tầng nghiệp vụ (services) tới tầng máy chủ HTTP (infra) và giao
diện người dùng (UI) cho cả Admin lẫn Customer.

\subsection{Ánh xạ khái niệm $\to$ file code $\to$ vai trò}

Bảng dưới đây liên kết từng khái niệm lý thuyết với file mã nguồn thực tế và vai
trò của nó trong hệ thống. Đây là bản đồ tổng quát để định vị nhanh khi vấn đáp.

\begin{table}[h]
\centering
\footnotesize
\begin{tabularx}{\linewidth}{L{3.3cm} L{4.6cm} X}
\toprule
\textbf{Khái niệm} & \textbf{File code} & \textbf{Vai trò} \\
\midrule
Vòng đời cert (status động) &
\code{services/cert\_lifecycle.py} &
Suy ra \code{active}/\code{expired}/\code{revoked}, liệt kê cert, revoke, renew. \\
\addlinespace
Workflow thu hồi 2 phía &
\code{services/revocation\_workflow.py} &
Customer gửi yêu cầu thu hồi; Admin duyệt (approve/reject) atomic. \\
\addlinespace
Lõi CRL \& OCSP DB &
\code{core/crl.py} &
Build/ký CRL bằng cryptography; đọc-ghi OCSP DB (JSON realtime). \\
\addlinespace
Snapshot \& publish CRL &
\code{services/crl\_publish.py} &
Chụp serial revoked từ DB, ký CRL bằng Root CA, ghi file + sync OCSP. \\
\addlinespace
Máy chủ phân phối CRL &
\code{infra/crl\_server.py} &
HTTP \code{GET /crl.pem} để client tải CRL về đối chiếu. \\
\addlinespace
OCSP responder &
\code{infra/ocsp\_server.py} &
HTTP \code{GET /ocsp?serial=} trả GOOD/REVOKED; flag \code{enabled} giả lập down. \\
\addlinespace
Quản lý vòng đời 4 server &
\code{services/infra\_manager.py} &
Singleton start/stop idempotent cặp Prod (8889/8888) + Lab (9889/9888). \\
\addlinespace
UI quản lý cert (Admin) &
\code{ui/admin/cert\_mgmt\_view.py} &
Bảng cert, dialog Revoke + Renew. \\
\addlinespace
UI duyệt thu hồi (Admin) &
\code{ui/admin/revoke\_queue\_view.py} &
Queue request, Approve (auto publish CRL) / Reject. \\
\addlinespace
UI cập nhật CRL (Admin) &
\code{ui/admin/crl\_publish\_view.py} &
Nút ``Publish CRL Now'', xem diff DB vs CRL hiện hành. \\
\addlinespace
UI cert của tôi (Customer) &
\code{ui/customer/my\_certs\_view.py} &
Customer xem + tải cert đã cấp. \\
\addlinespace
UI xin thu hồi (Customer) &
\code{ui/customer/revoke\_request\_view.py} &
Customer chọn cert active + lý do → gửi yêu cầu. \\
\addlinespace
UI tra CRL (Customer) &
\code{ui/customer/view\_crl\_view.py} &
Tra cứu danh sách thu hồi công khai theo serial/CN. \\
\bottomrule
\end{tabularx}
\end{table}

\begin{notebox}
\textbf{Phạm vi đáp ứng yêu cầu đề bài.} Cụm này hiện thực các hạng mục
A.8 (Admin quản lý cert: revoke/renew), A.9 (Admin duyệt yêu cầu thu hồi),
A.10 (cập nhật CRL), B.6 (Customer xem/tải cert), B.7 (Customer xin thu hồi),
B.8 (Customer tra CRL công khai) và một phần B.9 (hạ tầng để verify dựa vào
CRL/OCSP).
\end{notebox}

%==============================================================================
\section{Kiến trúc tổng thể và vị trí của thành viên trong hệ thống}
%==============================================================================

Hệ thống chia thành 3 tầng rõ rệt: tầng \textbf{lõi mật mã} (\code{core/}) chỉ
làm việc với đối tượng \code{cryptography}; tầng \textbf{nghiệp vụ}
(\code{services/}) chứa logic quy trình + truy vấn DB SQLite; tầng \textbf{hạ
tầng + UI} (\code{infra/}, \code{ui/}) cung cấp HTTP server và màn hình tương
tác. Cụm của Kiệt trải khắp cả 3 tầng theo một luồng dữ liệu thống nhất.

\subsection{Luồng thu hồi end-to-end}

Sơ đồ dưới đây mô tả luồng từ lúc khách hàng xin thu hồi cho tới lúc CRL/OCSP
phản ánh trạng thái mới ra bên ngoài.

\begin{center}
\flowbox{noteblue}{\textbf{Customer} chọn cert active + lý do\\
(\code{revoke\_request\_view} $\to$ \code{submit\_revoke\_request})}\\[4pt]
$\downarrow$ tạo bản ghi \code{revocation\_requests} status=\code{pending}\\[6pt]
\flowbox{warnbg}{\textbf{Admin} mở queue, bấm Approve\\
(\code{revoke\_queue\_view} $\to$ \code{approve\_revocation})}\\[4pt]
$\downarrow$ trong 1 transaction: set request=\code{approved} + UPDATE \code{issued\_certs.revoked\_at}\\[6pt]
\flowbox{okbg}{\code{sync\_ocsp\_db} ghi \code{ocsp\_db.json} (realtime)\\
$+$ \code{publish\_crl} ký CRL bằng Root CA $\to$ \code{crl.pem}}\\[4pt]
$\downarrow$\\[6pt]
\flowbox{codebg}{Hạ tầng phục vụ trạng thái mới:\\
\code{crl\_server} GET \code{/crl.pem} $\;|\;$ \code{ocsp\_server} GET \code{/ocsp?serial=}}
\end{center}

Điểm mấu chốt cần nhớ khi vấn đáp: có \textbf{hai nguồn sự thật tách biệt}.
\code{ocsp\_db.json} luôn được cập nhật \emph{ngay} mỗi khi có revoke (qua
\code{sync\_ocsp\_db}), nên OCSP responder luôn fresh. Còn file \code{crl.pem}
chỉ thay đổi khi có hành động \emph{Publish CRL} (thủ công hoặc tự động sau
approve). Sự lệch pha này chính là mô phỏng đúng thực tế: CRL có ``cửa sổ trễ''
giữa hai lần phát hành, còn OCSP gần như tức thời.

\subsection{Bảng vị trí module trong toàn hệ thống}

\begin{table}[h]
\centering
\footnotesize
\begin{tabularx}{\linewidth}{L{3.0cm} L{3.3cm} X}
\toprule
\textbf{Tầng} & \textbf{Module của Kiệt} & \textbf{Phụ thuộc / được gọi bởi} \\
\midrule
Core & \code{core/crl.py} & dùng \code{cryptography}; được \code{crl\_publish} và \code{ocsp\_server} gọi. \\
\addlinespace
Services & \code{cert\_lifecycle}, \code{revocation\_workflow}, \code{crl\_publish} & gọi \code{core/crl}, \code{ca\_admin} (load Root CA), \code{db.connection}. \\
\addlinespace
Services & \code{infra\_manager} & gọi \code{infra/crl\_server} + \code{infra/ocsp\_server}; được \code{ui/app.py} gọi lúc boot. \\
\addlinespace
Infra & \code{crl\_server}, \code{ocsp\_server} & dùng \code{http.server}; \code{ocsp\_server} đọc OCSP DB qua \code{core.crl}. \\
\addlinespace
UI Admin & 3 view & gọi services lifecycle/workflow/publish + \code{services.audit}. \\
\addlinespace
UI Customer & 3 view & gọi services hoặc Admin API (chế độ remote). \\
\bottomrule
\end{tabularx}
\end{table}

%==============================================================================
\section{Phần code cần làm (chi tiết)}
%==============================================================================

Phần này đi qua từng nhóm hàm quan trọng, kèm code thật trích từ mã nguồn và
giải thích từng ý. Mỗi sinh viên phụ trách cần đọc kỹ để trả lời được câu hỏi
``hàm này làm gì, vì sao viết như vậy''.

\subsection{Vòng đời cert: suy ra trạng thái động}

Điểm thiết kế cốt lõi của \code{cert\_lifecycle.py} là \textbf{status không lưu
trong DB} mà được \emph{tính lại mỗi lần đọc} dựa trên hai cột thô:
\code{revoked\_at} và \code{not\_valid\_after}.

\begin{Verbatim}
def _compute_status(row, now: datetime) -> str:
    """Suy ra status hiện tại của cert."""
    if row["revoked_at"]:
        return "revoked"
    na = _parse_iso(row["not_valid_after"])
    if na is None:
        return "active"
    if na.tzinfo is None:                # not_valid_after lưu tz-aware;
        na = na.replace(tzinfo=timezone.utc)   # nếu không, gán UTC.
    if na < now:
        return "expired"
    return "active"
\end{Verbatim}

\textbf{Giải thích từng ý:}
\begin{itemize}
  \item Thứ tự kiểm tra có chủ đích: \code{revoked} được ưu tiên cao nhất ---
    một cert đã thu hồi thì dù còn hạn vẫn báo \code{revoked}.
  \item Nếu chưa revoke, so sánh \code{not\_valid\_after} với \code{now}: quá
    hạn $\to$ \code{expired}, ngược lại $\to$ \code{active}.
  \item Xử lý timezone tỉ mỉ: nếu chuỗi ISO không kèm offset thì gán UTC để
    tránh lỗi so sánh ``naive vs aware datetime''.
\end{itemize}

\begin{notebox}
\textbf{Vì sao tính động thay vì lưu cột \code{status}?} Vì \code{expired} thay
đổi theo thời gian mà \emph{không} có sự kiện nào ghi vào DB --- một cert tự
``hết hạn'' khi đồng hồ vượt mốc. Nếu lưu cứng, ta phải có job quét định kỳ cập
nhật, dễ sai lệch. Tính động đảm bảo status luôn đúng tại thời điểm truy vấn.
\end{notebox}

Hàm \code{\_list\_certs} áp dụng filter status \emph{post-hoc} (sau khi đã tính
status) chứ không đưa vào câu \code{WHERE} SQL, đúng vì status là trường suy ra:

\begin{Verbatim}
now = datetime.now(timezone.utc)
with conn_scope(db_path) as conn:
    rows = conn.execute(sql, params).fetchall()
    out = []
    for r in rows:
        d = _row_to_dict(r, now)        # gắn d["status"]
        if status in (None, "all") or d["status"] == status:
            out.append(d)
    return out
\end{Verbatim}

\subsection{Revoke và Renew một chứng chỉ}

\code{revoke\_cert} đánh dấu thu hồi: validate lý do, chặn revoke trùng (giữ duy
nhất một sự kiện thu hồi), ghi \code{revoked\_at} trong transaction rồi
\emph{sync ngay} OCSP DB.

\begin{Verbatim}
now = datetime.now(timezone.utc).isoformat()
with transaction(db_path) as conn:
    row = conn.execute(
        "SELECT revoked_at FROM issued_certs WHERE id = ?", (cert_id,),
    ).fetchone()
    if row is None:
        raise CertLifecycleError("Không tìm thấy cert.")
    if row["revoked_at"]:
        raise CertLifecycleError(f"Cert đã bị revoked lúc {row['revoked_at']}.")
    conn.execute(
        "UPDATE issued_certs SET revoked_at = ?, revocation_reason = ? "
        "WHERE id = ?", (now, reason, cert_id),
    )
sync_ocsp_db(db_path, ocsp_db_path)     # OCSP fresh ngay, CRL publish sau
return get_cert_detail(cert_id, db_path)
\end{Verbatim}

\code{renew\_cert} \emph{phát hành cert mới} giữ nguyên subject + public key của
cert cũ, hiệu lực mới, không tự revoke cert cũ. Đây là khác biệt then chốt với
``reissue qua CSR mới'':

\begin{Verbatim}
old = get_cert_detail(cert_id, db_path)
if old["revoked_at"]:
    raise CertLifecycleError("Cert đã bị thu hồi — không thể renew.")
ca_cert, ca_key = load_active_root_ca_with_key(db_path)
old_cert_obj = x509.load_pem_x509_certificate(bytes(old["cert_pem"]))
new_cert, new_serial = reissue_cert_for_renewal(
    old_cert_obj, ca_cert, ca_key,
    validity_days=validity_days, ocsp_url=ocsp_url, crl_url=crl_url,
)
# INSERT cert mới với renewed_from_id = cert_id (truy vết chuỗi gia hạn)
\end{Verbatim}

\begin{notebox}
\textbf{Cột \code{renewed\_from\_id}} lưu liên kết cert mới $\to$ cert cũ, cho
phép truy vết toàn bộ ``chuỗi gia hạn'' của một domain. UI hiển thị cột ``Renew
từ'' chính là dựa vào trường này.
\end{notebox}

\subsection{Workflow thu hồi hai phía}

\code{revocation\_workflow.py} tách quyền: \textbf{customer chỉ được xin}, còn
\textbf{admin mới được quyết}. Phía customer, \code{submit\_revoke\_request}
chứa một \emph{BOLA guard} (Broken Object Level Authorization) quan trọng:

\begin{Verbatim}
cert = conn.execute(
    "SELECT id, owner_id, revoked_at, common_name "
    "FROM issued_certs WHERE id = ?", (cert_id,),
).fetchone()
if cert is None or cert["owner_id"] != requester_id:
    raise RevocationWorkflowError(
        f"Cert id={cert_id} không tồn tại hoặc không thuộc về bạn.")
if cert["revoked_at"]:
    raise RevocationWorkflowError("Cert đã bị thu hồi rồi, không cần gửi yêu cầu.")
# Chặn trùng: 1 cert chỉ có tối đa 1 request pending
dup = conn.execute(
    "SELECT id FROM revocation_requests "
    "WHERE issued_cert_id = ? AND status = 'pending'", (cert_id,),
).fetchone()
if dup:
    raise RevocationWorkflowError(f"Đã có yêu cầu pending (request #{dup['id']}).")
\end{Verbatim}

\textbf{Giải thích:} kiểm tra \code{owner\_id != requester\_id} ngăn người dùng A
xin thu hồi cert của người dùng B (lỗ hổng phân quyền cấp đối tượng cổ điển).
Việc chặn request trùng giữ queue admin sạch.

Phía admin, \code{approve\_revocation} là điểm \emph{atomic} chống race khi hai
admin cùng duyệt một request. Toàn bộ thay đổi DB nằm trong một transaction:

\begin{Verbatim}
with transaction(db_path) as conn:
    req = conn.execute("SELECT id, issued_cert_id, reason, status "
        "FROM revocation_requests WHERE id = ?", (req_id,)).fetchone()
    if req is None:
        raise RevocationWorkflowError("Không tìm thấy request.")
    if req["status"] != "pending":
        raise RevocationWorkflowError(
            f"Request đang ở status '{req['status']}', không approve được.")
    cert = conn.execute("SELECT id, revoked_at FROM issued_certs WHERE id = ?",
        (req["issued_cert_id"],)).fetchone()
    if cert["revoked_at"] is None:          # chỉ ghi nếu chưa revoke
        conn.execute("UPDATE issued_certs SET revoked_at = ?, "
            "revocation_reason = ? WHERE id = ?", (now, req["reason"], cert["id"]))
        cert_was_revoked = True
    else:
        cert_was_revoked = False            # admin đã revoke trực tiếp trước đó
    conn.execute("UPDATE revocation_requests SET status = 'approved', "
        "reviewed_at = ?, reviewed_by = ? WHERE id = ?", (now, admin_id, req_id))
\end{Verbatim}

Sau khi commit transaction, hàm gọi \code{sync\_ocsp\_db} (OCSP fresh ngay) và
nếu caller truyền \code{crl\_path} thì gọi luôn \code{publish\_crl} để CRL tự
cập nhật:

\begin{Verbatim}
sync_ocsp_db(db_path, ocsp_db_path)
crl_result = crl_error = None
if crl_path:
    try:
        crl_result = publish_crl(admin_id=admin_id, db_path=db_path,
            crl_path=crl_path, ocsp_db_path=ocsp_db_path)
    except CRLPublishError as e:
        crl_error = str(e)
return {"id": req_id, "cert_was_revoked": cert_was_revoked,
        "crl_result": crl_result, "crl_error": crl_error, ...}
\end{Verbatim}

\begin{notebox}
\textbf{Điểm tinh tế về race re-check.} Nếu admin đã revoke cert trực tiếp ở màn
A.8 trước đó, khi approve request hệ thống \emph{vẫn} đổi request sang
\code{approved} (để khép audit) nhưng \emph{không ghi đè} \code{revoked\_at} cũ
--- giữ đúng thời điểm thu hồi đầu tiên. Đây là re-check trạng thái cert ngay
trong transaction.
\end{notebox}

\subsection{Lõi CRL và OCSP DB (core/crl.py)}

\code{build\_crl} dựng đối tượng CRL bằng builder của thư viện
\code{cryptography} và \emph{ký} bằng khoá riêng của issuer (Root CA):

\begin{Verbatim}
def build_crl(issuer_cert, issuer_key, revoked_serials, validity_days=7):
    now = datetime.now(timezone.utc)
    builder = (x509.CertificateRevocationListBuilder()
        .issuer_name(issuer_cert.subject)
        .last_update(now)
        .next_update(now + timedelta(days=validity_days)))
    for serial in revoked_serials:
        revoked = (x509.RevokedCertificateBuilder()
            .serial_number(int(serial)).revocation_date(now).build())
        builder = builder.add_revoked_certificate(revoked)
    return builder.sign(private_key=issuer_key, algorithm=hashes.SHA256())
\end{Verbatim}

OCSP DB là một file JSON đơn giản chứa danh sách serial bị thu hồi. Hai hàm
đọc-ghi và một hàm dùng riêng để \emph{demo khoảng cách CRL vs OCSP}:

\begin{Verbatim}
def load_revoked_list(path: str) -> set:
    """Đọc danh sách serial bị revoke từ JSON. Trả về set[int]."""
    if not os.path.exists(path):
        return set()
    with open(path) as f:
        data = json.load(f)
    return set(int(s) for s in data)

def revoke_serial_ocsp_only(serial: int, ocsp_db_path: str):
    """Thêm serial vào OCSP DB mà KHÔNG cập nhật CRL."""
    revoked = load_revoked_list(ocsp_db_path)
    revoked.add(serial)
    save_revoked_list(list(revoked), ocsp_db_path)
\end{Verbatim}

\code{build\_and\_publish\_crl} là tiện ích snapshot toàn bộ OCSP DB rồi ký CRL
--- tương đương bấm nút ``Publish CRL Now''.

\subsection{Snapshot và publish CRL (crl\_publish.py)}

\code{snapshot\_revoked\_serials} đọc tất cả cert có \code{revoked\_at} khác
NULL, parse \code{serial\_hex} (hex) sang int. Đáng chú ý: \emph{không} loại trừ
cert đã hết hạn --- đúng chuẩn CRL.

\begin{Verbatim}
def snapshot_revoked_serials(db_path: str) -> "list[int]":
    with conn_scope(db_path) as conn:
        rows = conn.execute(
            "SELECT serial_hex FROM issued_certs "
            "WHERE revoked_at IS NOT NULL ORDER BY revoked_at ASC").fetchall()
    out = []
    for r in rows:
        try:
            out.append(int(r["serial_hex"], 16))
        except (ValueError, TypeError):
            print(f"[crl_publish] WARN: serial_hex không parse được: "
                  f"{r['serial_hex']!r}", file=sys.stderr)
    return out
\end{Verbatim}

\code{publish\_crl} là trái tim của A.10: load Root CA active, snapshot serial,
build + ký CRL, ghi file, đồng bộ OCSP, trả metadata cho UI.

\begin{Verbatim}
def publish_crl(admin_id, db_path, crl_path=DEFAULT_CRL_PATH,
                ocsp_db_path=DEFAULT_OCSP_DB_PATH, validity_days=7):
    try:
        ca_cert, ca_key = load_active_root_ca_with_key(db_path)
    except CAError as e:
        raise CRLPublishError(f"Không publish được CRL: {e} Tạo Root CA trước.")
    serials = snapshot_revoked_serials(db_path)
    crl = build_crl(ca_cert, ca_key, serials, validity_days=validity_days)
    os.makedirs(os.path.dirname(crl_path) or ".", exist_ok=True)
    save_crl(crl, crl_path)
    if ocsp_db_path:
        sync_ocsp_db(db_path, ocsp_db_path)
    return {"revoked_count": len(serials), "crl_path": os.path.abspath(crl_path),
            "this_update": ..., "next_update": ...,
            "issuer": ca_cert.subject.rfc4514_string()}
\end{Verbatim}

\code{sync\_ocsp\_db} tách riêng vì \emph{không cần ký} --- chỉ ghi JSON, nên
revoke flow gọi được tức thì mà không phụ thuộc Root CA. \code{list\_crl\_entries}
và \code{get\_published\_crl\_info} parse lại file CRL để UI hiển thị; bản trước
còn JOIN \code{issued\_certs} làm giàu thêm CN + owner cho dễ đọc.

\subsection{Hạ tầng máy chủ: CRL server}

\code{crl\_server.py} dùng factory pattern: mỗi lần gọi
\code{start\_crl\_server} tạo một Handler class riêng capture \code{crl\_path}
qua closure, cho phép chạy nhiều instance (prod + lab) trên port khác nhau.

\begin{Verbatim}
class CRLHandler(BaseHTTPRequestHandler):
    def do_GET(self):
        if self.path not in ("/crl.pem", "/crl"):
            self.send_response(404); self.end_headers(); return
        if not os.path.exists(crl_path):
            self.send_response(404); self.end_headers()
            self.wfile.write(b"CRL not found"); return
        with open(crl_path, "rb") as f:
            data = f.read()
        self.send_response(200)
        self.send_header("Content-Type", "application/x-pem-file")
        self.send_header("Content-Length", str(len(data)))
        self.end_headers(); self.wfile.write(data)
\end{Verbatim}

Server chạy trên \code{ThreadingHTTPServer} với \code{daemon\_threads = True}:
mỗi request một thread con, một client chậm không treo các client khác, và
thread tự thoát khi process chính kết thúc.

\subsection{Hạ tầng máy chủ: OCSP responder}

\code{ocsp\_server.py} nhận \code{GET /ocsp?serial=<int>} và trả JSON
\code{GOOD}/\code{REVOKED}. Flag \code{enabled} bọc trong một dict mutable để UI
\emph{toggle runtime} giả lập OCSP down (trả 503) mà không cần restart server.

\begin{Verbatim}
class OCSPHandler(BaseHTTPRequestHandler):
    def do_GET(self):
        if not state.get("enabled", True):           # giả lập down
            body = b'{"error": "OCSP responder temporarily unavailable"}'
            self.send_response(503); ...; self.wfile.write(body); return
        parsed = urlparse(self.path)
        if parsed.path != "/ocsp":
            self.send_response(404); self.end_headers(); return
        serial_raw = parse_qs(parsed.query).get("serial", [None])[0]
        if serial_raw is None:
            self._json(400, {"error": "missing 'serial' parameter"}); return
        serial_int = int(serial_raw)
        revoked = load_revoked_list(revoked_list_path)   # đọc JSON realtime
        status = "REVOKED" if serial_int in revoked else "GOOD"
        self._json(200, {"serial": str(serial_int), "status": status})
\end{Verbatim}

\begin{notebox}
\textbf{Vì sao OCSP đọc file mỗi request?} Để luôn fresh. Mỗi khi revoke gọi
\code{sync\_ocsp\_db} ghi lại JSON, lần tra OCSP kế tiếp đã thấy ngay. Đây chính
là lý do OCSP ``realtime'' còn CRL thì trễ.
\end{notebox}

\subsection{Quản lý vòng đời 4 server: infra\_manager}

\code{infra\_manager.py} là singleton quản hai cặp server: \textbf{Prod}
(8889/8888, auto-start lúc boot, phục vụ cert thật từ DB) và \textbf{Lab}
(9889/9888, start thủ công khi mở Verification Lab, phục vụ cert demo). Các thao
tác start đều \emph{idempotent} --- gọi nhiều lần chỉ start một lần:

\begin{Verbatim}
def start_prod_servers(self):
    os.makedirs(CERTS_DIR, exist_ok=True)
    if self._prod_crl is None:                  # idempotent guard
        try:
            self._prod_crl = start_crl_server(host="localhost",
                port=PROD_CRL_PORT, crl_path=PROD_CRL_PATH)
        except OSError as e:
            _log.warning("Prod CRL server không start được (port %d?): %s",
                         PROD_CRL_PORT, e)
    if self._prod_ocsp is None:
        self._prod_ocsp, self._prod_ocsp_state = start_ocsp_server(
            host="localhost", port=PROD_OCSP_PORT, revoked_list_path=PROD_OCSP_DB)
\end{Verbatim}

Toggle OCSP enabled dựa vào việc \code{start\_ocsp\_server} trả về cả
\emph{state dict mutable}; manager giữ reference này và sửa trực tiếp field
\code{enabled}:

\begin{Verbatim}
def set_lab_ocsp_enabled(self, enabled: bool):
    if self._lab_ocsp_state is not None:
        self._lab_ocsp_state["enabled"] = enabled   # handler thấy ngay
\end{Verbatim}

Singleton lấy qua \code{get\_infra()} (lazy init), đảm bảo cả app chỉ có một
instance quản lý vòng đời 4 server suốt process.

\subsection{Giao diện người dùng (UI)}

\textbf{Admin --- \code{cert\_mgmt\_view}}: bảng cert tô màu theo status
(active xanh, expired xám, revoked đỏ), nút Revoke mở \code{RevokeCertDialog}
(lý do bắt buộc) gọi \code{revoke\_cert}, nút Renew mở \code{RenewCertDialog} gọi
\code{renew\_cert}; mọi thao tác đều ghi \code{write\_audit}.

\textbf{Admin --- \code{revoke\_queue\_view}}: queue request (mặc định filter
\code{pending}). Khi Approve, gọi \code{approve\_revocation} với cả
\code{crl\_path} nên CRL \emph{tự publish}; UI ghi audit \code{REVOKE\_APPROVED},
\code{CERT\_REVOKED} và \code{CRL\_PUBLISHED}, và thông báo số serial trong CRL.

\begin{Verbatim}
result = approve_revocation(req_id, self.app.session["id"], self.app.db_path,
    ocsp_db_path=DEFAULT_OCSP_DB_PATH, crl_path=DEFAULT_CRL_PATH)
...
crl_result = result.get("crl_result")
if crl_result:
    write_audit(..., Action.CRL_PUBLISHED, ...,
        details={"revoked_count": crl_result["revoked_count"],
                 "source": "auto_after_revoke_approve"})
\end{Verbatim}

\textbf{Admin --- \code{crl\_publish\_view}}: hiển thị metadata CRL hiện hành,
snapshot DB (số serial sẽ vào CRL), \emph{diff} cảnh báo nếu DB lệch CRL, và nút
``Publish CRL Now'' gọi \code{publish\_crl}.

\textbf{Customer --- \code{my\_certs\_view}}: xem + tải cert của mình
(\code{list\_certs\_for\_owner}, có chế độ remote qua Admin API).
\textbf{\code{revoke\_request\_view}}: chọn cert active + lý do, gọi
\code{submit\_revoke\_request}. \textbf{\code{view\_crl\_view}}: tra CRL công
khai theo serial/CN qua \code{list\_crl\_entries} + \code{get\_published\_crl\_info}.

%==============================================================================
\section{Cần nắm để vấn đáp}
%==============================================================================

\begin{qabox}
\textbf{Hỏi:} Trạng thái cert được lưu cứng trong DB hay suy ra động? Vì sao?\\[2pt]
\textbf{Đáp:} Suy ra động trong \code{\_compute\_status}. DB chỉ lưu hai cột thô
\code{revoked\_at} và \code{not\_valid\_after}; status \code{active}/
\code{expired}/\code{revoked} được tính lại mỗi lần đọc. Lý do: \code{expired}
phụ thuộc thời gian hiện tại, không có ``sự kiện hết hạn'' để ghi DB --- tính
động luôn đúng, tránh dữ liệu lỗi thời.
\end{qabox}

\begin{qabox}
\textbf{Hỏi:} CRL và OCSP khác nhau ra sao và bù nhau thế nào?\\[2pt]
\textbf{Đáp:} CRL là \emph{danh sách} toàn bộ serial bị thu hồi do CA ký, phát
hành định kỳ (\code{this\_update} $\to$ \code{next\_update}); client tải về một
lần rồi tra offline --- ưu điểm là không cần online liên tục, nhược điểm là có
độ trễ tới lần publish kế tiếp và file phình to. OCSP tra \emph{từng serial}
theo thời gian thực, trả GOOD/REVOKED ngay --- ưu điểm tươi mới, nhược điểm là
phụ thuộc responder online. Trong đồ án, OCSP DB cập nhật tức thì mỗi lần revoke
(\code{sync\_ocsp\_db}), còn CRL chỉ đổi khi Publish --- minh hoạ đúng sự bù trừ
giữa hai cơ chế.
\end{qabox}

\begin{qabox}
\textbf{Hỏi:} Vì sao CRL phải được ký bằng Root CA?\\[2pt]
\textbf{Đáp:} CRL là tuyên bố ``những cert này không còn tin cậy''. Nếu không
ký, kẻ tấn công có thể giả mạo CRL rỗng để giấu việc một cert đã bị thu hồi,
hoặc ngược lại đưa cert hợp lệ vào CRL để từ chối dịch vụ. Chữ ký của Root CA
(\code{builder.sign(issuer\_key, SHA256)}) cho phép client xác minh CRL thật sự
do CA phát hành và chưa bị sửa. OCSP DB JSON \emph{không ký} chỉ vì nó là kênh
nội bộ phục vụ responder, không phát tán trực tiếp như CRL.
\end{qabox}

\begin{qabox}
\textbf{Hỏi:} Mô tả luồng thu hồi hai phía (customer xin, admin duyệt).\\[2pt]
\textbf{Đáp:} Customer gọi \code{submit\_revoke\_request} (có BOLA guard kiểm tra
sở hữu, chặn cert đã revoked, chặn request pending trùng), tạo bản ghi
\code{pending}. Admin mở queue, Approve gọi \code{approve\_revocation}: trong một
transaction set request=\code{approved} và UPDATE \code{revoked\_at}, sau commit
thì sync OCSP + publish CRL. Nếu cert đã revoke trước đó thì giữ nguyên thời
điểm cũ. Admin cũng có thể Reject (lý do bắt buộc, append vào reason để giữ
history).
\end{qabox}

\begin{qabox}
\textbf{Hỏi:} \code{renew\_cert} khác ``reissue qua CSR mới'' ở điểm nào?\\[2pt]
\textbf{Đáp:} Renew giữ nguyên \emph{subject + public key} của cert cũ, chỉ đổi
serial + hiệu lực mới, ghi \code{renewed\_from\_id} để truy vết --- customer
không phải gửi CSR lại. Reissue qua CSR là cấp một cert hoàn toàn mới (có thể
khoá mới). Renew từ chối nếu cert cũ đã thu hồi, và \emph{không} tự revoke cert
cũ (admin tự quyết).
\end{qabox}

\begin{qabox}
\textbf{Hỏi:} Vòng đời các server hạ tầng được quản lý ra sao?\\[2pt]
\textbf{Đáp:} \code{InfraManager} là singleton (lấy qua \code{get\_infra()}).
Prod (8889/8888) auto-start lúc boot và chạy suốt đời app; Lab (9889/9888) start
thủ công khi mở Verification Lab, stop khi đóng. Start idempotent nhờ guard
\code{if self.\_prod\_crl is None}. \code{stop\_all} gọi shutdown + server\_close
cả 4. OCSP có thể bật/tắt runtime qua state dict mutable.
\end{qabox}

\begin{qabox}
\textbf{Hỏi:} Vì sao tách \code{ocsp\_db.json} và \code{crl.pem} thành hai nguồn?
\\[2pt]
\textbf{Đáp:} Để mô phỏng đúng thực tế PKI. \code{ocsp\_db.json} cập nhật tức thì
mỗi revoke (không cần ký) nên OCSP luôn fresh; \code{crl.pem} chỉ đổi khi Publish
(cần ký Root CA) nên có cửa sổ trễ. Màn \code{crl\_publish\_view} hiển thị diff
giữa số serial trong DB và trong CRL để nhắc admin publish khi lệch.
\end{qabox}

\begin{warnbox}
\textbf{Lưu ý dễ bị hỏi vặn.} \code{serial\_hex} trong DB là chuỗi hex, phải
\code{int(serial\_hex, 16)} mới ra số để đưa vào CRL/OCSP. OCSP responder so
sánh \code{serial\_int in revoked} với \code{revoked} là \code{set[int]} ---
nếu lưu sai kiểu (str vs int) sẽ luôn trả GOOD nhầm. Hãy giải thích được điểm
này.
\end{warnbox}

%==============================================================================
\section{Test case và demo}
%==============================================================================

\subsection{Bảng test case chính}

\begin{table}[h]
\centering
\footnotesize
\begin{tabularx}{\linewidth}{L{3.4cm} X L{3.1cm}}
\toprule
\textbf{Test} & \textbf{Kịch bản} & \textbf{Kỳ vọng} \\
\midrule
\code{compute\_status active} & cert chưa revoke, còn hạn & status = \code{active} \\
\addlinespace
\code{compute\_status expired} & \code{not\_valid\_after} < now & status = \code{expired} \\
\addlinespace
\code{compute\_status revoked} & \code{revoked\_at} != NULL & status = \code{revoked} (ưu tiên) \\
\addlinespace
\code{revoke\_cert twice} & revoke 1 cert hai lần & lần 2 raise \code{CertLifecycleError} \\
\addlinespace
\code{renew\_cert} & renew cert active & cert mới, \code{renewed\_from\_id} đúng \\
\addlinespace
\code{renew revoked} & renew cert đã revoke & raise lỗi \\
\addlinespace
\code{submit BOLA} & customer xin revoke cert người khác & raise (không thuộc về bạn) \\
\addlinespace
\code{submit duplicate} & xin revoke khi đã có pending & raise (đã có request) \\
\addlinespace
\code{approve atomic} & approve request pending & request=approved + cert revoked \\
\addlinespace
\code{approve twice} & approve request đã approved & raise (status không pending) \\
\addlinespace
\code{publish\_crl} & có 2 cert revoked & CRL ký, \code{revoked\_count}=2 \\
\addlinespace
\code{OCSP GOOD/REVOKED} & tra serial chưa/đã revoke & trả GOOD / REVOKED \\
\addlinespace
\code{OCSP disabled} & set enabled=False rồi tra & HTTP 503 \\
\addlinespace
\code{infra idempotent} & gọi \code{start\_prod} 2 lần & chỉ 1 server, không lỗi \\
\bottomrule
\end{tabularx}
\end{table}

\subsection{Đoạn unit test gợi ý}

Test vòng đời + thu hồi (gợi ý cho \code{test\_m7\_cert\_lifecycle.py} và
\code{test\_m8\_revocation\_crl.py}):

\begin{Verbatim}
def test_revoke_then_status_revoked(tmp_db):
    cert = issue_sample_cert(tmp_db)
    revoke_cert(cert["id"], admin_id=1, reason="key lộ", db_path=tmp_db)
    detail = get_cert_detail(cert["id"], tmp_db)
    assert detail["status"] == "revoked"
    assert detail["revoked_at"] is not None

def test_submit_bola_guard(tmp_db):
    cert = issue_sample_cert(tmp_db, owner_id=10)
    with pytest.raises(RevocationWorkflowError):
        submit_revoke_request(cert["id"], requester_id=99,
                              reason="x", db_path=tmp_db)  # không phải chủ

def test_publish_crl_signs_with_root_ca(tmp_db):
    cert = issue_sample_cert(tmp_db)
    revoke_cert(cert["id"], 1, "r", tmp_db)
    res = publish_crl(admin_id=1, db_path=tmp_db,
                      crl_path=tmp_path/"crl.pem")
    assert res["revoked_count"] == 1
    crl = load_crl(tmp_path/"crl.pem")
    assert crl.is_signature_valid(root_ca_public_key)
\end{Verbatim}

Test hạ tầng (gợi ý cho \code{test\_infra\_manager.py}):

\begin{Verbatim}
def test_start_prod_idempotent():
    mgr = InfraManager()
    mgr.start_prod_servers()
    first = mgr._prod_crl
    mgr.start_prod_servers()          # gọi lần 2
    assert mgr._prod_crl is first     # không tạo server mới
    mgr.stop_all()

def test_ocsp_toggle_returns_503(ocsp_running):
    server, state = ocsp_running
    state["enabled"] = False
    resp = http_get("http://localhost:8888/ocsp?serial=1")
    assert resp.status == 503
\end{Verbatim}

\subsection{Kịch bản demo cho hội đồng}

\begin{enumerate}
  \item \textbf{Revoke + publish:} Admin mở Quản lý chứng nhận, chọn một cert
    active, bấm Revoke (nhập lý do). Mở Cập nhật CRL: thấy diff ``DB có 1
    revoked, CRL có 0''. Bấm \emph{Publish CRL Now} $\to$ CRL ký lại,
    \code{revoked\_count}=1.
  \item \textbf{Tra CRL/OCSP:} Customer mở Tra cứu CRL thấy serial vừa revoke.
    Gọi \code{GET /ocsp?serial=<serial>} (trình duyệt hoặc curl) trả
    \code{REVOKED}; tra một serial active trả \code{GOOD}.
  \item \textbf{Luồng hai phía:} Customer mở Yêu cầu thu hồi, chọn cert active +
    lý do, Gửi. Admin mở Duyệt yêu cầu thu hồi, Approve $\to$ cert revoked + CRL
    tự publish (kiểm tra thông báo số serial).
  \item \textbf{Giả lập OCSP down:} Trong Verification Lab, tắt OCSP (toggle
    \code{enabled=False}); tra OCSP trả HTTP 503 --- minh hoạ vì sao cần CRL làm
    phương án dự phòng khi responder gặp sự cố.
\end{enumerate}

\begin{notebox}
\textbf{Lệnh chạy test.} Từ thư mục gốc dự án:
\code{pytest tests/test\_m7\_cert\_lifecycle.py tests/test\_m8\_revocation\_crl.py
tests/test\_infra\_manager.py -v}. Để demo OCSP bằng dòng lệnh:
\code{curl "http://localhost:8888/ocsp?serial=12345"}.
\end{notebox}

%==============================================================================
\section{Tích hợp vào chương trình chung}
%==============================================================================

Cụm này không đứng riêng mà ghép vào toàn hệ thống qua các điểm nối sau:

\begin{itemize}
  \item \textbf{Khởi động app:} \code{ui/app.py} gọi
    \code{get\_infra().start\_prod\_servers()} lúc boot để CRL + OCSP server
    Prod sẵn sàng phục vụ; khi thoát gọi \code{stop\_all()}.
  \item \textbf{Phụ thuộc Root CA:} \code{publish\_crl} và \code{renew\_cert}
    gọi \code{load\_active\_root\_ca\_with\_key} (module CA của thành viên khác).
    Nếu chưa có Root CA, hai hàm raise lỗi rõ ràng để UI hiển thị.
  \item \textbf{Nhúng URL vào cert:} khi cấp/renew cert, URL CRL + OCSP
    (\code{prod\_crl\_url()}, \code{prod\_ocsp\_url()} từ \code{infra\_manager})
    được nhúng vào extension CRL Distribution Points / AIA --- để khâu
    \emph{chain validation} (phần việc Phạm Minh Quân) biết đi tra ở đâu.
  \item \textbf{Audit log:} mọi hành động revoke/renew/approve/reject/publish đều
    gọi \code{services.audit.write\_audit} với \code{Action} tương ứng
    (\code{CERT\_REVOKED}, \code{CERT\_RENEWED}, \code{REVOKE\_APPROVED},
    \code{CRL\_PUBLISHED}\ldots) --- ghép vào màn Audit chung.
  \item \textbf{Chế độ remote:} các view Customer (\code{my\_certs},
    \code{revoke\_request}, \code{view\_crl}) hỗ trợ gọi Admin API
    (\code{remote\_csr\_client}) khi chạy tách máy Admin/Customer --- cùng một
    logic nghiệp vụ nhưng qua HTTP.
  \item \textbf{Cấu hình:} đường dẫn \code{PROD\_CRL\_PATH},
    \code{PROD\_OCSP\_DB\_PATH} và port lấy từ \code{config.py} (cho phép override
    qua biến môi trường để tránh xung đột port).
\end{itemize}

\begin{notebox}
\textbf{Quan hệ với chain validation.} Đầu ra của cụm này (CRL đã ký phục vụ tại
\code{/crl.pem} và OCSP responder tại \code{/ocsp}) chính là \emph{đầu vào} cho
bước kiểm tra trạng thái thu hồi trong quy trình verify của Phạm Minh Quân. Hai
phần ghép với nhau tạo thành vòng tin cậy đầy đủ: cấp $\to$ thu hồi $\to$ công bố
$\to$ kiểm tra.
\end{notebox}

%==============================================================================
\section{Bổ sung sau rà soát: các hàm/feature dễ bị hỏi}

\begin{itemize}
  \item \textbf{\code{unrevoke\_serial(serial, ocsp\_db\_path)} --- gỡ thu hồi.}
        Đối xứng với \code{revoke\_serial\_ocsp\_only}: loại serial khỏi OCSP DB
        (rollback, vd xóa server demo). Điểm nhấn: \emph{CRL KHÔNG tự đổi} --- phải
        \emph{Publish CRL Now} mới đồng bộ. Đây là chiều ngược của lệch pha CRL vs
        OCSP (OCSP cập nhật tức thì, CRL theo chu kỳ).
  \item \textbf{Hai hàm liệt kê thu hồi song song.} \code{list\_pending\_revocations}
        (hardcode \code{WHERE status='pending'}) là tiện ích queue thuần; nhưng UI
        \code{revoke\_queue\_view} thực tế gọi \code{list\_all\_revocations(status=...)}
        (mặc định pending) để combobox cho xem cả approved/rejected/all.
  \item \textbf{Các API truy vấn trạng thái của \code{InfraManager}.} Ngoài
        \code{set\_lab\_ocsp\_enabled}: \code{status()} trả dict 4 cờ để dashboard vẽ
        badge; \code{is\_prod\_running}/\code{is\_lab\_running} chỉ True khi \emph{cả}
        CRL$+$OCSP chạy; \code{get\_lab\_ocsp\_state} trả dict \emph{mutable} để Lab
        bind \code{BooleanVar} toggle trực tiếp; \code{set\_prod\_ocsp\_enabled} toggle
        Prod (chủ yếu cho test).
  \item \textbf{BOLA guard tầng 2 --- \code{get\_cert\_detail(cert\_id, owner\_id=...)}}.
        Khi customer xem chi tiết, truyền \code{owner\_id} $\Rightarrow$ thêm
        \code{AND c.owner\_id = ?} và trả \code{None} nếu cert không thuộc về họ ---
        chặn customer đoán \code{cert\_id} để xem/tải cert người khác (ngoài BOLA ở
        \code{submit\_revoke\_request}).
  \item \textbf{Cặp ghi-đọc serial trong \code{ocsp\_db.json}.} \code{save\_revoked\_list}
        cố ý ghi serial dạng \emph{chuỗi} (\code{str(s)} --- JSON không an toàn cho
        số nguyên rất lớn); \code{load\_revoked\_list} ép \code{int(s)} khi đọc. Cặp
        đối xứng này là lý do so sánh \code{serial in revoked} (set \code{int}) luôn
        đúng kiểu --- khép kín cảnh báo ``sai kiểu str/int luôn trả GOOD nhầm''.
\end{itemize}

%==============================================================================
\section{Checklist chuẩn bị vấn đáp}
%==============================================================================

\begin{itemize}
  \item[\textbf{$\square$}] Giải thích được vì sao status cert tính động, đọc
    trôi chảy \code{\_compute\_status} và thứ tự ưu tiên revoked/expired/active.
  \item[\textbf{$\square$}] Phân biệt rõ CRL vs OCSP: cơ chế, ưu/nhược, cách bù
    nhau; chỉ ra trong code đâu là realtime (\code{sync\_ocsp\_db}) đâu là trễ
    (\code{publish\_crl}).
  \item[\textbf{$\square$}] Nói được vì sao CRL phải ký bằng Root CA và OCSP DB
    thì không cần ký.
  \item[\textbf{$\square$}] Vẽ lại luồng thu hồi hai phía và chỉ ra điểm atomic +
    BOLA guard trong code.
  \item[\textbf{$\square$}] Phân biệt \code{renew\_cert} (giữ subject + public
    key) với reissue qua CSR mới; giải thích \code{renewed\_from\_id}.
  \item[\textbf{$\square$}] Trình bày vòng đời 4 server, tính idempotent của
    start, và cơ chế toggle OCSP qua state dict mutable.
  \item[\textbf{$\square$}] Chạy được demo: revoke $\to$ publish $\to$ tra
    CRL/OCSP $\to$ tắt OCSP giả lập down (HTTP 503).
  \item[\textbf{$\square$}] Chạy được bộ test \code{test\_m7\_cert\_lifecycle},
    \code{test\_m8\_revocation\_crl}, \code{test\_infra\_manager} và giải thích
    từng assert.
  \item[\textbf{$\square$}] Giải thích điểm nối với chain validation và lý do
    nhúng URL CRL/OCSP vào cert lúc phát hành.
  \item[\textbf{$\square$}] Trả lời được vì sao \code{serial\_hex} cần
    \code{int(.., 16)} và rủi ro sai kiểu str/int trong so sánh OCSP.
\end{itemize}

\end{document}
